\scanpage{369}
\mspage{93}
$X$ et $Y$, j'entends leurs maquettes, sont bien homotopes dans $\mathrm{Cat}$, s'il est vrai qu'une subdivision simpliciale (il suffirait : la subdivision barycentrique, \inserted{en itérant\ldots}) \deleted{\unsure{de}} de $X$ \deleted{\unsure{est}}\unsure{a} une maquette dans $\mathrm{Cat}$ qui lui est homotope. La question soulevée ici me paraît intéressante, il faut visiblement en trouver la réponse. Elle se pose déjà pour \deleted{tout} ensemble ordonné $I$, qu'on va supposer fini pour simplifier, en prenant l'ensemble (ordonné \inserted{par inclusion}) $I'=\mathrm{sd}(I)$ des parties totalement ordonnées de $I$. On sait que les réalisations topologiques de $I$ et de $\mathrm{sd}(I)$ sont canoniquement \unsure{homéomorphes}, mais $I$ et $\mathrm{sd}(I)$ sont-ils homotopes dans $\mathrm{Cat}$ ? On a bien une application croissante

\[
 \text{can.}\qquad f:\mathrm{sd}(I)\longrightarrow I,
\]

qui associe à toute partie tot. ordonnée de $I$ son plus grand élément. Mais existe-t-il une \unsure{application} croissante en sens inverse, qui soit un homotopisme inverse du précédent à homotopie près ?

\deleted{\unsure{Je}} Je viens de vérifier dans le cas
% archive sheet number: 368
